% #############################################################################
% Abstract Text
% !TEX root = ../main.tex
% #############################################################################
% reset acronyms
\acresetall

\begingroup
\singlespacing
\setlength{\parskip}{8pt}
\setlength{\parindent}{1.5em}

% use \noindent in firts paragraph
% \noindent
\noindent European Union regulation imposes organisational obligations whose implications for individual stakeholders are often difficult to identify. Translating these obligations into representations of daily activities, responsibilities, and dependencies requires a systematic connection between legal requirements and operational evidence. This dissertation proposes a human-supervised method, supported by \ac{GenAI}, for identifying, specifying, and materialising stakeholder-specific architecture viewpoints.

Following \ac{DSRM} and grounded in ISO/IEC/IEEE 42010, the method combines stakeholder interviews, operational-footprint extraction, regulatory relevance analysis, viewpoint specification, and ArchiMate model generation. Traceability links architectural elements and relationships to binding legal text and documented operational evidence, while explicit ownership distinctions, human approval gates, and cross-artifact reconciliation constrain generated content.

The method is demonstrated using the \ac{DORA} through the perspectives of impacted stakeholders. The resulting \ac{AD}s relate regulatory requirements to role-specific work while distinguishing direct responsibilities, external dependencies, and uncertain ownership. Preliminary stakeholder feedback indicates that the formal views can support regulatory understanding and reveal missing operational context, enabling viewpoint refinement. It also identifies limitations in explanatory narratives and simplified diagrams. The demonstrations support the feasibility of the approach, with opportunities for further development. The contribution is a traceable process for communicating stakeholder-specific regulatory impact without assessing organisational compliance.

\endgroup